% GENERATED FILE. Edit canonical lessons, then run npm run generate:books.
% canonical-source-hash: fnv1a64:a9b3af56576141e9
% canonical-lessons: RU-C402-podklyuchat, RU-C402-zaryazhat, RU-C402-perepisyvatsya, RU-C402-peresylat, RU-C402-ustroystvo, RU-R402-review-401-402, RU-R402-review-final-402

\chapter{Staying connected}
\label{ch:ru-staying-connected}
\hlchaptermodality{402}

\begin{chapteropening}
\textbf{By the end of this chapter:} \emph{I can connect a device, charge my phone, message a friend and forward a letter.}

Complete the last lesson of chapter 402: I can connect a device, charge my phone, message a friend and forward a letter.
\end{chapteropening}

\section[\texorpdfstring{\ru{подключать} (podklyuchát)}{подключать (podklyuchát)}]{\ru{подключать} (podklyuchát) — to connect, to plug in}

\label{lesson:RU-C402-podklyuchat}



\begin{quote}
Before the new one: say the Russian for to delete, then the Russian for to load.
\end{quote}

\subsection*{You'll want to know: \ru{подключать}}
\textbf{\ru{подключать}} (\emph{podklyuchát}) — \textquotedblleft{}to connect, to plug in\textquotedblright{}.

\begin{quote}
\ru{Мы} \ru{подключаем} \ru{интернет} \ru{в} \ru{новой} \ru{квартире}.
\end{quote}

\begin{quote}
\emph{My podklyucháyem internét v nóvoy kvartíre.}
\end{quote}

\begin{quote}
We are getting the internet connected in the new flat.
\end{quote}

Imperfective; perfective \emph{\ru{подключить}}. \emph{\ru{Подключаться} \ru{к}} + dative — to connect to.

\begin{cousinweb}
\textbf{\ru{подключать}} is built from \emph{\ru{под}-} and \emph{\ru{ключ}} \textquotedblleft{}a key\textquotedblright{}, inherited from Proto-Slavic \emph{*ključĭ}: to lock one thing into another.
\end{cousinweb}

\begin{grammarlens}[title={What you've built}]
One more word for telling, explaining and describing at B1.
\end{grammarlens}

\subsection*{Your turn}
\begin{itemize}
\raggedright
  \item \emph{Say it:} \emph{podklyuchát}
  \item \emph{Say it:} \emph{podklyuchát}, once more
  \item \emph{Say it:} a sentence of your own with \emph{podklyuchát}
  \item \emph{Recall:} say the Russian for to delete, then the Russian for to load, then say \emph{podklyuchát} again
\end{itemize}

\begin{tcolorbox}[breakable,colback=teal!4,colframe=teal!35!black,title={Before you move on}]
What does \ru{подключать} mean? (\textquotedblleft{}To connect, to plug in\textquotedblright{}.) Say it once more.
\end{tcolorbox}
\section[\texorpdfstring{\ru{заряжать} (zaryazhát)}{заряжать (zaryazhát)}]{\ru{заряжать} (zaryazhát) — to charge (a battery)}

\label{lesson:RU-C402-zaryazhat}



\begin{quote}
Before the new one: say the Russian for to load, then the Russian for to connect.
\end{quote}

\subsection*{You'll want to know: \ru{заряжать}}
\textbf{\ru{заряжать}} (\emph{zaryazhát}) — \textquotedblleft{}to charge (a battery)\textquotedblright{}.

\begin{quote}
\ru{Я} \ru{заряжаю} \ru{телефон} \ru{каждую} \ru{ночь}.
\end{quote}

\begin{quote}
\emph{Ya zaryazháyu telefón kázhduyu noch.}
\end{quote}

\begin{quote}
I charge my phone every night.
\end{quote}

Imperfective; perfective \emph{\ru{зарядить}}. \emph{\ru{Зарядка}} is a charger, and also morning exercises.

\begin{cousinweb}
\textbf{\ru{заряжать}} is built from \emph{\ru{за}-} and the root of \emph{\ru{ряд}} \textquotedblleft{}a row, an order\textquotedblright{}, as in \emph{\ru{снаряжать}} \textquotedblleft{}to equip\textquotedblright{}: first to load a gun, now to charge a battery.
\end{cousinweb}

\begin{grammarlens}[title={What you've built}]
One more word for telling, explaining and describing at B1.
\end{grammarlens}

\subsection*{Your turn}
\begin{itemize}
\raggedright
  \item \emph{Say it:} \emph{zaryazhát}
  \item \emph{Say it:} \emph{zaryazhát}, once more
  \item \emph{Say it:} a sentence of your own with \emph{zaryazhát}
  \item \emph{Recall:} say the Russian for to load, then the Russian for to connect, then say \emph{zaryazhát} again
\end{itemize}

\begin{tcolorbox}[breakable,colback=teal!4,colframe=teal!35!black,title={Before you move on}]
What does \ru{заряжать} mean? (\textquotedblleft{}To charge (a battery)\textquotedblright{}.) Say it once more.
\end{tcolorbox}
\section[\texorpdfstring{\ru{переписываться} (perepísyvatsya)}{переписываться (perepísyvatsya)}]{\ru{переписываться} (perepísyvatsya) — to correspond, to message each other}

\label{lesson:RU-C402-perepisyvatsya}



\begin{quote}
Before the new one: say the Russian for to connect, then the Russian for to charge.
\end{quote}

\subsection*{You'll want to know: \ru{переписываться}}
\textbf{\ru{переписываться}} (\emph{perepísyvatsya}) — \textquotedblleft{}to correspond, to message each other\textquotedblright{}.

\begin{quote}
\ru{Мы} \ru{переписываемся} \ru{с} \ru{другом} \ru{уже} \ru{пять} \ru{лет}.
\end{quote}

\begin{quote}
\emph{My perepísyvayemsya s drúgom uzhé pyat let.}
\end{quote}

\begin{quote}
My friend and I have been corresponding for five years.
\end{quote}

Imperfective; perfective \emph{\ru{переписаться}}. \emph{\ru{Переписываться} \ru{с}} + instrumental. Without \emph{-\ru{ся}}, \emph{\ru{переписывать}} means to copy out.

\begin{cousinweb}
\textbf{\ru{переписываться}} is built from \emph{\ru{пере}-} \textquotedblleft{}across, back and forth\textquotedblright{} and \emph{\ru{писать}} \textquotedblleft{}to write\textquotedblright{}, inherited from Proto-Slavic \emph{*pisati} \textquotedblleft{}to write; to paint\textquotedblright{}, from PIE \emph{*pey\'{k}-} \textquotedblleft{}to paint, to colour\textquotedblright{} (EtymDB): letters going back and forth.
\end{cousinweb}

\begin{grammarlens}[title={What you've built}]
One more word for telling, explaining and describing at B1.
\end{grammarlens}

\subsection*{Your turn}
\begin{itemize}
\raggedright
  \item \emph{Say it:} \emph{perepísyvatsya}
  \item \emph{Say it:} \emph{perepísyvatsya}, once more
  \item \emph{Say it:} a sentence of your own with \emph{perepísyvatsya}
  \item \emph{Recall:} say the Russian for to connect, then the Russian for to charge, then say \emph{perepísyvatsya} again
\end{itemize}

\begin{tcolorbox}[breakable,colback=teal!4,colframe=teal!35!black,title={Before you move on}]
What does \ru{переписываться} mean? (\textquotedblleft{}To correspond, to message each other\textquotedblright{}.) Say it once more.
\end{tcolorbox}
\section[\texorpdfstring{\ru{пересылать} (peresylát)}{пересылать (peresylát)}]{\ru{пересылать} (peresylát) — to forward, to send on}

\label{lesson:RU-C402-peresylat}



\begin{quote}
Before the new one: say the Russian for to charge, then the Russian for to correspond.
\end{quote}

\subsection*{You'll want to know: \ru{пересылать}}
\textbf{\ru{пересылать}} (\emph{peresylát}) — \textquotedblleft{}to forward, to send on\textquotedblright{}.

\begin{quote}
\ru{Я} \ru{пересылаю} \ru{тебе} \ru{его} \ru{письмо}.
\end{quote}

\begin{quote}
\emph{Ya peresyláyu tebé yegó pismó.}
\end{quote}

\begin{quote}
I am forwarding his letter to you.
\end{quote}

Imperfective; perfective \emph{\ru{переслать}} (\emph{\ru{перешлю}, \ru{перешлёшь}}).

\begin{cousinweb}
\textbf{\ru{пересылать}} is built from \emph{\ru{пере}-} \textquotedblleft{}across, on\textquotedblright{} and \emph{\ru{слать}} \textquotedblleft{}to send\textquotedblright{}, the root of \emph{\ru{посылка}}: to send a thing on.
\end{cousinweb}

\begin{grammarlens}[title={What you've built}]
One more word for telling, explaining and describing at B1.
\end{grammarlens}

\subsection*{Your turn}
\begin{itemize}
\raggedright
  \item \emph{Say it:} \emph{peresylát}
  \item \emph{Say it:} \emph{peresylát}, once more
  \item \emph{Say it:} a sentence of your own with \emph{peresylát}
  \item \emph{Recall:} say the Russian for to charge, then the Russian for to correspond, then say \emph{peresylát} again
\end{itemize}

\begin{tcolorbox}[breakable,colback=teal!4,colframe=teal!35!black,title={Before you move on}]
What does \ru{пересылать} mean? (\textquotedblleft{}To forward, to send on\textquotedblright{}.) Say it once more.
\end{tcolorbox}
\section[\texorpdfstring{\ru{устройство} (ustróystvo)}{устройство (ustróystvo)}]{\ru{устройство} (ustróystvo) — a device; the structure, the set-up}

\label{lesson:RU-C402-ustroystvo}



\begin{quote}
Before the new one: say the Russian for to correspond, then the Russian for to forward.
\end{quote}

\subsection*{You'll want to know: \ru{устройство}}
\textbf{\ru{устройство}} (\emph{ustróystvo}) — \textquotedblleft{}a device; the structure, the set-up\textquotedblright{}.

\begin{quote}
\ru{Это} \ru{устройство} \ru{работает} \ru{без} \ru{интернета}.
\end{quote}

\begin{quote}
\emph{Éto ustróystvo rabótayet bez internéta.}
\end{quote}

\begin{quote}
This device works without the internet.
\end{quote}

Neuter. \emph{\ru{Мобильное} \ru{устройство}} — a mobile device. Also the way a thing is arranged: \emph{\ru{устройство} \ru{общества}}.

\begin{cousinweb}
\textbf{\ru{устройство}} is built on \emph{\ru{устроить}} \textquotedblleft{}to arrange, to set up\textquotedblright{}, from \emph{\ru{строить}} \textquotedblleft{}to build\textquotedblright{}, on the old Slavic root \emph{*stroji} \textquotedblleft{}order, arrangement\textquotedblright{}.
\end{cousinweb}

\begin{grammarlens}[title={What you've built}]
One more word for telling, explaining and describing at B1.
\end{grammarlens}

\subsection*{Your turn}
\begin{itemize}
\raggedright
  \item \emph{Say it:} \emph{ustróystvo}
  \item \emph{Say it:} \emph{ustróystvo}, once more
  \item \emph{Say it:} a sentence of your own with \emph{ustróystvo}
  \item \emph{Recall:} say the Russian for to correspond, then the Russian for to forward, then say \emph{ustróystvo} again
\end{itemize}

\begin{tcolorbox}[breakable,colback=teal!4,colframe=teal!35!black,title={Before you move on}]
What does \ru{устройство} mean? (\textquotedblleft{}A device; the structure, the set-up\textquotedblright{}.) Say it once more.
\end{tcolorbox}
\section[(dialogue)]{Conditions and rules — a review}

\label{lesson:RU-R402-review-401-402}



\begin{quote}
Say the Russian for an app and to install.
\end{quote}

\subsection*{Your turn}
Say these aloud:

\begin{itemize}
\raggedright
  \item an app, to install, to update, to delete, to load
  \item to connect, to charge, to correspond, to forward, a device
\end{itemize}

\begin{tcolorbox}[breakable,colback=teal!4,colframe=teal!35!black,title={Before you move on}]
An app? (\textbf{\ru{приложение}}.) To install? (\textbf{\ru{устанавливать}}.) To update? (\textbf{\ru{обновлять}}.) To delete? (\textbf{\ru{удалять}}.) To load? (\textbf{\ru{загружать}}.) To connect? (\textbf{\ru{подключать}}.) To charge? (\textbf{\ru{заряжать}}.) To correspond? (\textbf{\ru{переписываться}}.) To forward? (\textbf{\ru{пересылать}}.) A device? (\textbf{\ru{устройство}}.)
\end{tcolorbox}
\section[(dialogue)]{Staying connected — once more}

\label{lesson:RU-R402-review-final-402}



\begin{quote}
Say the Russian for to connect and to charge.
\end{quote}

\subsection*{Your turn}
\begin{itemize}
\raggedright
  \item \emph{Say it:} to connect, to charge, to correspond, to forward, a device
\end{itemize}

\begin{tcolorbox}[breakable,colback=teal!4,colframe=teal!35!black,title={Before you move on}]
To connect? (\textbf{\ru{подключать}}.) To charge? (\textbf{\ru{заряжать}}.) To correspond? (\textbf{\ru{переписываться}}.) To forward? (\textbf{\ru{пересылать}}.) A device? (\textbf{\ru{устройство}}.)
\end{tcolorbox}
